\documentclass[10pt, aspectratio=169]{beamer}
\usetheme{Madrid}
\usecolortheme{default}

\usepackage[utf8]{inputenc}
\usepackage[T1]{fontenc}
\usepackage[portuguese]{babel}
\usepackage{xcolor}

\title[Auditoria Vanna AI]{Auditoria Forense e Plano de Resgate Técnico}
\subtitle{Ecossistema Vanna AI}

\author[Equipe de Auditoria]{
    Bruno Henrique Carneiro da Silva \and Diego Carvalho Cavalcante \\
    Fábio Henrique Lisboa de Souza \and Flávio Henrique de Jesus Cruz \\
    Guilherme Sollano Andrade dos Santos \and José Weverton de Oliveira Vilar \\
    Leonardo Caricchio do Nascimento \and Matheus Henrique Silva de Melo
}

\institute[COMP0503]{Engenharia de Software (COMP0503)}
\date{02 de Junho de 2026}

\begin{document}

\begin{frame}
    \titlepage
\end{frame}

\begin{frame}{Avaliação de Gestão (MPS.BR - GPR)}
    \begin{itemize}
        \item \textbf{Risco: Lacuna na Engenharia de Requisitos (Mudança de Escopo oculta)}
        \begin{itemize}
            \item \textcolor{blue}{\textbf{Evidência:}} Issue \#659 e PR \#660. Uma discussão iniciada como correção de erro (\textit{documentation\_collection}) escalou para refatoração de \textit{embeddings} sem nenhum RFC ou documento de arquitetura formalizando a alteração.
        \end{itemize}
        \vspace{0.3cm}
        
        \item \textbf{Risco: Vulnerabilidade de Manutenção ("Herói Único")}
        \begin{itemize}
            \item \textcolor{blue}{\textbf{Evidência:}} Gráfico de contribuições do repositório revela picos severos de \textit{crunch time} e um mantenedor central responsável por 399 \textit{commits}, contra apenas 60 do segundo contribuidor mais ativo.
        \end{itemize}
    \end{itemize}
\end{frame}

\begin{frame}{Anatomia do Código e Dívida Técnica (SOLID e DRY)}
    \begin{itemize}
        \item \textbf{Risco: Violação Crítica do Princípio da Responsabilidade Única (SRP)}
        \begin{itemize}
            \item \textcolor{blue}{\textbf{Evidência:}} Arquivo \texttt{src/vanna/legacy/base/base.py}. Trata-se de uma \textit{God Class} confirmada com exatas 2.125 linhas e 65 métodos, aglutinando geração de SQL, manipulação de banco de dados e execução de IA no mesmo escopo.
        \end{itemize}
        \vspace{0.3cm}

        \item \textbf{Risco: Retrabalho e Quebra Parcial do DRY}
        \begin{itemize}
            \item \textcolor{blue}{\textbf{Evidência:}} O método \texttt{\_extract\_tool\_calls\_from\_message} possui a mesma lógica de tratamento de JSON repetida em diferentes implementações de \texttt{LlmService} (OpenAI, Anthropic), em vez de estar isolada em uma classe utilitária.
        \end{itemize}
    \end{itemize}
\end{frame}

\begin{frame}{Riscos em Padrões de Projeto (GoF)}
    \begin{itemize}
        \item \textbf{Risco: Falha de Gerenciamento Criacional (Conexões Redundantes)}
        \begin{itemize}
            \item \textcolor{blue}{\textbf{Evidência:}} O arquivo \texttt{chromadb/agent\_memory.py} invoca \texttt{chromadb.PersistentClient()} a cada nova instância, sem a proteção de um padrão \textit{Singleton} para reaproveitar a conexão vetorial.
        \end{itemize}
        \vspace{0.3cm}

        \item \textbf{Risco: Delegação Insegura de Execução (Ausência de Padrão Comportamental)}
        \begin{itemize}
            \item \textcolor{blue}{\textbf{Evidência:}} O \textit{Agentic Loop} (\texttt{core/agent/agent.py}) confia cegamente no modelo. Ao identificar o comando, executa \texttt{await self.tool\_registry.execute(tool\_call, context)} sem nenhuma corrente de validação prévia.
        \end{itemize}
    \end{itemize}
\end{frame}

\begin{frame}{Plano de Resgate: Falha Arquitetural}
    \textbf{Diagnóstico da Vulnerabilidade no Vanna 2.0}
    \vspace{0.3cm}
    \begin{itemize}
        \item \textbf{Assimetria de Design:} O projeto utiliza de forma inteligente a abstração de \textit{Middlewares} para a comunicação externa com as APIs de LLM, mas ignora completamente este padrão na sua camada mais sensível: a execução de ferramentas.
        \item \textbf{Execução Direta:} O método \texttt{await self.tool\_registry.execute()} roda a ação baseada unicamente na saída textual gerada pela Inteligência Artificial.
        \item \textbf{Risco Crítico:} Uma eventual alucinação do modelo contendo instruções destrutivas (ex: \texttt{DROP} ou \texttt{DELETE}) atinge o banco de dados de produção instantaneamente, sem nenhum mecanismo de contenção.
    \end{itemize}
\end{frame}

\begin{frame}{Plano de Resgate: \textit{Chain of Responsibility} (GoF)}
    \textbf{Proposta: Escudo de \textit{Middlewares} de Execução}
    \vspace{0.3cm}
    \begin{itemize}
        \item \textbf{Isolamento de Responsabilidades (SRP):} Criação da entidade \texttt{ToolExecutionRequest} para circular por uma corrente de interceptadores independentes.
        \item \textbf{Elo 1 - \texttt{RateLimitMiddleware}:} Validação das cotas de uso do usuário, prevenindo gargalos operacionais ou custos excessivos de API.
        \item \textbf{Elo 2 - \texttt{SQLSanitizationMiddleware}:} Inspeção rigorosa dos argumentos da chamada (ex: \texttt{args.sql}) para barrar operações DDL/DML não autorizadas.
        \item \textbf{Elo 3 - \texttt{FinalExecutionMiddleware}:} Camada de execução física (\texttt{sql\_runner}), atingida exclusivamente se a requisição sobreviver intacta a todos os elos anteriores.
    \end{itemize}
\end{frame}

\begin{frame}{Roadmap de Maturidade: Rumo ao Nível G (MPS.BR)}
    \textbf{Objetivo:} Institucionalizar processos em 90 dias para superar o atual estágio informal (Pré-Nível G).
    \vspace{0.3cm}

    \textbf{Ação 1: Institucionalizar a Gerência de Projetos (GPR)}
    \begin{itemize}
        \item \textbf{Foco:} Dissolver o anti-padrão de "Herói Único" e garantir cadência previsível.
        \item \textbf{Implementação Tática:} 
        \begin{itemize}
            \item Implantação de \textit{GitHub Projects} (\textit{Backlog, In Progress, Review, Done}).
            \item Política de rodízio: designar pelo menos dois \textit{co-maintainers} com permissões de \textit{merge}.
            \item Criação de um \textbf{Registro de Riscos} formal no \textit{wiki} (mapeando instabilidade de APIs e \textit{vendor lock-in}).
        \end{itemize}
        \item \textbf{Prazo:} Semanas 1 a 3 \textcolor{red}{\textbf{(Prioridade Crítica)}}.
    \end{itemize}
\end{frame}

\begin{frame}{Roadmap de Maturidade: Requisitos e Qualidade}
    \textbf{Ação 2: Formalizar a Gerência de Requisitos (GRE)}
    \begin{itemize}
        \item \textbf{Foco:} Erradicar mudanças de escopo ocultas e garantir rastreabilidade bidirecional.
        \item \textbf{Implementação Tática:} \textit{Issue Templates} obrigatórios (com Critérios de Aceitação) e uso de \textit{ADRs} (\textit{Architecture Decision Records}) para documentar alterações estruturais.
        \item \textbf{Prazo:} Semanas 2 a 5 \textcolor{orange}{\textbf{(Prioridade Alta)}}.
    \end{itemize}
    \vspace{0.3cm}

    \textbf{Ação 3: Elevar a Qualidade do Produto (QPR)}
    \begin{itemize}
        \item \textbf{Foco:} Proteger a base de código contra regressões e vulnerabilidades de IA.
        \item \textbf{Implementação Tática:} Estabelecer \textit{Definition of Done} (DoD) inflexível. \textit{Pull Requests} bloqueados nativamente até aprovação em CI/CD (cobertura \textit{pytest} $\ge$ 70\%) e varredura SAST (Bandit) contra injeções SQL.
        \item \textbf{Prazo:} Semanas 3 a 8 \textcolor{orange}{\textbf{(Prioridade Alta)}}.
    \end{itemize}
\end{frame}

\end{document}